% GENERATED FILE. Edit canonical lessons, then run npm run generate:books.
% canonical-source-hash: fnv1a64:2318290a2323d847
% canonical-lessons: GE-C311-weihnachtsmarkt, GE-C311-weihnachtsbaum, GE-C311-adventskranz, GE-C311-bescherung, GE-C311-lebkuchen

\chapter{Im Advent}
\label{ch:ge-im-advent}
\hlchaptermodality{311}

\begin{chapteropening}
\textbf{By the end of this chapter:} \emph{I can talk about Christmas customs: a Christmas market, a Christmas tree, an Advent wreath, the giving of presents and gingerbread.}

Complete the last lesson of chapter 311: I can talk about Christmas customs: a Christmas market, a Christmas tree, an Advent wreath, the giving of presents and gingerbread.
\end{chapteropening}

\section[der Weihnachtsmarkt]{der Weihnachtsmarkt — the Christmas market}

\label{lesson:GE-C311-weihnachtsmarkt}



\begin{quote}
Before the new one: say the German for to dress up, then the German for the custom.
\end{quote}

\subsection*{You'll want to know: der Weihnachtsmarkt}
\textbf{der Weihnachtsmarkt} — \textquotedblleft{}the Christmas market\textquotedblright{}.

\begin{quote}
Auf dem Weihnachtsmarkt tranken wir heißen Tee.
\end{quote}

\begin{quote}
\emph{At the Christmas market we drank hot tea.}
\end{quote}

Plural: \emph{die Weihnachtsmärkte}.

\begin{cousinweb}
\emph{Weihnachtsmarkt} joins \emph{Weihnachten}, Middle High German \emph{wīhenahten} \textquotedblleft{}the holy nights\textquotedblright{} (EtymDB), and \emph{Markt}, which came from Latin \emph{mercātus} \textquotedblleft{}market\textquotedblright{} (EtymDB).
\end{cousinweb}

\begin{grammarlens}[title={What you've built}]
One more word for telling, explaining and describing at B1.
\end{grammarlens}

\subsection*{Your turn}
\begin{itemize}
\raggedright
  \item \emph{Say it:} \emph{der Weihnachtsmarkt}
  \item \emph{Say it:} \emph{der Weihnachtsmarkt}, once more
  \item \emph{Say it:} a sentence of your own with \emph{der Weihnachtsmarkt}
  \item \emph{Recall:} say the German for to dress up, then the German for the custom, then say \emph{der Weihnachtsmarkt} again
\end{itemize}

\begin{tcolorbox}[breakable,colback=teal!4,colframe=teal!35!black,title={Before you move on}]
What does der Weihnachtsmarkt mean? (\textquotedblleft{}The Christmas market\textquotedblright{}.) Say it once more.
\end{tcolorbox}
\section[der Weihnachtsbaum]{der Weihnachtsbaum — the Christmas tree}

\label{lesson:GE-C311-weihnachtsbaum}



\begin{quote}
Before the new one: say the German for the custom, then the German for the Christmas market.
\end{quote}

\subsection*{You'll want to know: der Weihnachtsbaum}
\textbf{der Weihnachtsbaum} — \textquotedblleft{}the Christmas tree\textquotedblright{}.

\begin{quote}
Der Weihnachtsbaum stand im Wohnzimmer.
\end{quote}

\begin{quote}
\emph{The Christmas tree stood in the living room.}
\end{quote}

Plural: \emph{die Weihnachtsbäume}. In the south and in Austria also \emph{der Christbaum}.

\begin{cousinweb}
\emph{Weihnachtsbaum} joins \emph{Weihnachten} \textquotedblleft{}Christmas\textquotedblright{} and \emph{Baum} \textquotedblleft{}tree\textquotedblright{}, Proto-Germanic \emph{baumaz}, a cousin of English \emph{beam} (EtymDB).
\end{cousinweb}

\begin{grammarlens}[title={What you've built}]
One more word for telling, explaining and describing at B1.
\end{grammarlens}

\subsection*{Your turn}
\begin{itemize}
\raggedright
  \item \emph{Say it:} \emph{der Weihnachtsbaum}
  \item \emph{Say it:} \emph{der Weihnachtsbaum}, once more
  \item \emph{Say it:} a sentence of your own with \emph{der Weihnachtsbaum}
  \item \emph{Recall:} say the German for the custom, then the German for the Christmas market, then say \emph{der Weihnachtsbaum} again
\end{itemize}

\begin{tcolorbox}[breakable,colback=teal!4,colframe=teal!35!black,title={Before you move on}]
What does der Weihnachtsbaum mean? (\textquotedblleft{}The Christmas tree\textquotedblright{}.) Say it once more.
\end{tcolorbox}
\section[der Adventskranz]{der Adventskranz — the Advent wreath}

\label{lesson:GE-C311-adventskranz}



\begin{quote}
Before the new one: say the German for the Christmas market, then the German for the Christmas tree.
\end{quote}

\subsection*{You'll want to know: der Adventskranz}
\textbf{der Adventskranz} — \textquotedblleft{}the Advent wreath\textquotedblright{}.

\begin{quote}
Auf dem Adventskranz brannten zwei Kerzen.
\end{quote}

\begin{quote}
\emph{Two candles burned on the Advent wreath.}
\end{quote}

Plural: \emph{die Adventskränze}. Four candles, one more each Sunday.

\begin{cousinweb}
\emph{Adventskranz} joins \emph{Advent}, from Latin \emph{adventus} \textquotedblleft{}arrival\textquotedblright{}, and \emph{Kranz} \textquotedblleft{}wreath\textquotedblright{}, Old High German \emph{kranz} (EtymDB).
\end{cousinweb}

\begin{grammarlens}[title={What you've built}]
One more word for telling, explaining and describing at B1.
\end{grammarlens}

\subsection*{Your turn}
\begin{itemize}
\raggedright
  \item \emph{Say it:} \emph{der Adventskranz}
  \item \emph{Say it:} \emph{der Adventskranz}, once more
  \item \emph{Say it:} a sentence of your own with \emph{der Adventskranz}
  \item \emph{Recall:} say the German for the Christmas market, then the German for the Christmas tree, then say \emph{der Adventskranz} again
\end{itemize}

\begin{tcolorbox}[breakable,colback=teal!4,colframe=teal!35!black,title={Before you move on}]
What does der Adventskranz mean? (\textquotedblleft{}The Advent wreath\textquotedblright{}.) Say it once more.
\end{tcolorbox}
\section[die Bescherung]{die Bescherung — the giving of Christmas presents}

\label{lesson:GE-C311-bescherung}



\begin{quote}
Before the new one: say the German for the Christmas tree, then the German for the Advent wreath.
\end{quote}

\subsection*{You'll want to know: die Bescherung}
\textbf{die Bescherung} — \textquotedblleft{}the giving of Christmas presents\textquotedblright{}.

\begin{quote}
Die Bescherung war am Abend.
\end{quote}

\begin{quote}
\emph{The presents were given out in the evening.}
\end{quote}

In most families on the evening of the twenty-fourth.

\begin{cousinweb}
\emph{Bescherung} is built on \emph{bescheren} \textquotedblleft{}to give, to allot\textquotedblright{}, from Middle High German \emph{beschern} \textquotedblleft{}to destine, to allot\textquotedblright{} (EtymDB).
\end{cousinweb}

\begin{grammarlens}[title={What you've built}]
One more word for telling, explaining and describing at B1.
\end{grammarlens}

\subsection*{Your turn}
\begin{itemize}
\raggedright
  \item \emph{Say it:} \emph{die Bescherung}
  \item \emph{Say it:} \emph{die Bescherung}, once more
  \item \emph{Say it:} a sentence of your own with \emph{die Bescherung}
  \item \emph{Recall:} say the German for the Christmas tree, then the German for the Advent wreath, then say \emph{die Bescherung} again
\end{itemize}

\begin{tcolorbox}[breakable,colback=teal!4,colframe=teal!35!black,title={Before you move on}]
What does die Bescherung mean? (\textquotedblleft{}The giving of Christmas presents\textquotedblright{}.) Say it once more.
\end{tcolorbox}
\section[der Lebkuchen]{der Lebkuchen — the gingerbread}

\label{lesson:GE-C311-lebkuchen}



\begin{quote}
Before the new one: say the German for the Advent wreath, then the German for the giving of Christmas presents.
\end{quote}

\subsection*{You'll want to know: der Lebkuchen}
\textbf{der Lebkuchen} — \textquotedblleft{}the gingerbread\textquotedblright{}.

\begin{quote}
Meine Oma backte jedes Jahr Lebkuchen.
\end{quote}

\begin{quote}
\emph{My grandma baked gingerbread every year.}
\end{quote}

Plural: \emph{die Lebkuchen}.

\begin{cousinweb}
\emph{Lebkuchen} comes from Middle High German \emph{lebekuoche} (EtymDB), with \emph{Kuchen} \textquotedblleft{}cake\textquotedblright{}; the origin of the first part is not settled.
\end{cousinweb}

\begin{grammarlens}[title={What you've built}]
One more word for telling, explaining and describing at B1.
\end{grammarlens}

\subsection*{Your turn}
\begin{itemize}
\raggedright
  \item \emph{Say it:} \emph{der Lebkuchen}
  \item \emph{Say it:} \emph{der Lebkuchen}, once more
  \item \emph{Say it:} a sentence of your own with \emph{der Lebkuchen}
  \item \emph{Recall:} say the German for the Advent wreath, then the German for the giving of Christmas presents, then say \emph{der Lebkuchen} again
\end{itemize}

\begin{tcolorbox}[breakable,colback=teal!4,colframe=teal!35!black,title={Before you move on}]
What does der Lebkuchen mean? (\textquotedblleft{}The gingerbread\textquotedblright{}.) Say it once more.
\end{tcolorbox}
